% --- Couverture -------------------------------------------------------------
\begin{titlepage}
\thispagestyle{empty}
\begin{tikzpicture}[remember picture, overlay]

  % Fond papier
  \fill[bmtPapier] (current page.south west) rectangle (current page.north east);

  % Bloc principal indigo
  \fill[bmtIndigo]
    (current page.north west) rectangle
    ([yshift=-19.4cm]current page.north east);

  % --- Composition géométrique : un signe par domaine de l'année ------------
  % Analyse : une courbe de logarithme et sa tangente
  \begin{scope}[shift={([xshift=1.5cm, yshift=-9.1cm]current page.north west)}]
    \clip (-0.1,-0.15) rectangle (5.4,4.6);
    \draw[bmtPapier, line width=1.5pt, opacity=0.92, smooth, samples=160,
          domain=0.09:5.3]
      plot (\x, {1.35 + 1.5*ln(\x)});
    \draw[bmtRaphia, line width=1.1pt, dash pattern=on 3.5pt off 3pt]
      (0.35, 0) -- (4.6, 4.25);
  \end{scope}

  % Algèbre : le système linéaire 3x3
  \foreach \i in {0,1,2} {
    \foreach \j in {0,1,2} {
      \fill[bmtLaterite, opacity={0.35 + 0.2*\i}]
        ([xshift={8.6cm + \j*0.68cm}, yshift={-4.5cm - \i*0.68cm}]current page.north west)
        rectangle ++(0.48cm, -0.48cm);}}

  % Suites : les termes d'une suite géométrique décroissante
  \begin{scope}[shift={([xshift=11.9cm, yshift=-7.6cm]current page.north west)}]
    \draw[bmtPapier, line width=1pt, opacity=0.55] (-0.2,0) -- (3.5,0);
    \foreach \k [count=\i from 0] in {1.85, 1.24, 0.83, 0.55, 0.37, 0.25}{
      \draw[bmtPalissandre!85!bmtPapier, line width=1pt]
        ({\i*0.6},0) -- ({\i*0.6},\k);
      \fill[bmtPalissandre!70!bmtPapier] ({\i*0.6},\k) circle (0.085);}
  \end{scope}

  % Probabilités : un arbre à deux niveaux
  \begin{scope}[shift={([xshift=8.9cm, yshift=-10.4cm]current page.north west)}]
    \draw[bmtRaphia, line width=1pt] (0,0) -- (1.0,0.54);
    \draw[bmtRaphia, line width=1pt] (0,0) -- (1.0,-0.54);
    \draw[bmtRaphia, line width=1pt, opacity=0.65] (1.0,0.54) -- (2.0,0.88);
    \draw[bmtRaphia, line width=1pt, opacity=0.65] (1.0,0.54) -- (2.0,0.24);
    \draw[bmtRaphia, line width=1pt, opacity=0.65] (1.0,-0.54) -- (2.0,-0.24);
    \draw[bmtRaphia, line width=1pt, opacity=0.65] (1.0,-0.54) -- (2.0,-0.88);
    \foreach \p in {(0,0), (1.0,0.54), (1.0,-0.54)}
      \fill[bmtPapier] \p circle (0.075);
  \end{scope}

  % Statistique : un nuage de points et sa droite d'ajustement
  \begin{scope}[shift={([xshift=13.4cm, yshift=-11.4cm]current page.north west)}]
    \draw[bmtPapier, line width=1pt, opacity=0.5] (-0.15,0) -- (3.1,0);
    \draw[bmtPapier, line width=1pt, opacity=0.5] (0,-0.15) -- (0,2.1);
    \foreach \x/\y in {0.35/0.45, 0.75/0.72, 1.15/0.66, 1.6/1.15,
                       2.0/1.28, 2.45/1.62}
      \fill[bmtVetiver!75!bmtPapier] (\x,\y) circle (0.085);
    \draw[bmtRaphia, line width=1.1pt] (0.2,0.32) -- (2.75,1.78);
  \end{scope}

  % --- Titre ---------------------------------------------------------------
  \node[anchor=north west, text=bmtPapier, align=left,
        font=\sffamily\fontsize{12.5}{15}\selectfont, opacity=0.85]
    at ([xshift=1.95cm, yshift=-13.35cm]current page.north west)
    {BOKY MATEMATIKA};

  \node[anchor=north west, text=bmtPapier, align=left,
        font=\sffamily\bfseries\fontsize{54}{58}\selectfont]
    at ([xshift=1.8cm, yshift=-14.05cm]current page.north west)
    {Terminale~L};

  \draw[bmtRaphia, line width=3pt]
    ([xshift=1.95cm, yshift=-16.75cm]current page.north west) --
    ([xshift=5.65cm, yshift=-16.75cm]current page.north west);

  \node[anchor=north west, text=bmtPapier, align=left, text width=13.6cm,
        font=\sffamily\fontsize{10.5}{15}\selectfont, opacity=0.92]
    at ([xshift=1.95cm, yshift=-17.35cm]current page.north west)
    {Algèbre \quad Analyse \quad Suites numériques \quad Traitement des
     données et probabilités};

  % --- Bas de couverture ---------------------------------------------------
  \node[anchor=north west, text=bmtEncre, align=left, text width=12cm,
        font=\sffamily\fontsize{11}{15}\selectfont]
    at ([xshift=1.9cm, yshift=-20.6cm]current page.north west)
    {\textbf{Programme d'Études 2026}\\
     Enseignement Secondaire Général — République de Madagascar\\
     Épreuve de mathématiques générales, options A1 et A2};

  % Marque de l'éditeur
  \draw[bmtIndigo, line width=1.6pt]
    ([xshift=1.95cm, yshift=-22.95cm]current page.north west) --
    ([xshift=3.35cm, yshift=-22.95cm]current page.north west);
  \node[anchor=north west, text=bmtIndigo, align=left,
        font=\sffamily\bfseries\fontsize{12}{15}\selectfont]
    at ([xshift=1.9cm, yshift=-23.35cm]current page.north west)
    {ÉDITIONS MADA-BOKY};

  % Les quatre couleurs des quatre parties, en signature de bas de page
  \foreach \c [count=\i from 0] in {bmtLaterite, bmtIndigo, bmtPalissandre, bmtRaphia}{
    \fill[\c] ([xshift={1.9cm + \i*0.95cm}, yshift=2.5cm]current page.south west)
      rectangle ++(0.62cm, 0.62cm);}

  \node[anchor=south west, text=bmtGris, align=left,
        font=\sffamily\fontsize{9}{12}\selectfont]
    at ([xshift=1.9cm, yshift=1.5cm]current page.south west)
    {\ifbmtprof ÉDITION DU PROFESSEUR\else ÉDITION DE L'ÉLÈVE\fi};

\end{tikzpicture}
\end{titlepage}
